\section*{Chapitre \ref{ch:g10:stats} --- Statistique descriptive}

\begin{solution}{exo:g10:stats:1}
Somme~: $0+1+1+2+2+2+3+3+4+4+5+6+7+8+12 = 60$, donc la \omterm{def:g10:stats:mean}{moyenne} est
$\frac{60}{15} = 4$. La série est déjà triée et $N = 15$ est impair~: la
\omterm{def:g10:stats:median}{médiane} est la $8$-ième valeur, $3$. \omterm{def:g10:stats:quartiles}{Étendue}~: $12 - 0 = 12$.
\end{solution}

\begin{solution}{exo:g10:stats:2}
\omterm{def:g10:stats:series}{Fréquences}~: $\frac{7}{50} = 0.14$, $\frac{9}{50} = 0.18$,
$\frac{8}{50} = 0.16$, $\frac{10}{50} = 0.20$, $\frac{6}{50} = 0.12$,
$\frac{10}{50} = 0.20$ (somme $1$). Issue \omterm{def:g10:stats:mean}{moyenne}~:
\[
\bar x = \frac{7 \times 1 + 9 \times 2 + 8 \times 3 + 10 \times 4
+ 6 \times 5 + 10 \times 6}{50}
= \frac{7 + 18 + 24 + 40 + 30 + 60}{50}
= \frac{179}{50} = 3.58 .
\]
\end{solution}

\begin{solution}{exo:g10:stats:3}
$N = 12$ (pair)~: la \omterm{def:g10:stats:median}{médiane} est le \omterm{prop:g10:coordgeom:midpoint}{milieu} des $6$-ième et $7$-ième
valeurs, toutes deux égales à $6$~: \omterm{def:g10:stats:median}{médiane} $6$.

$Q_1$~: $\frac{12}{4} = 3$, donc $Q_1$ est la $3$-ième valeur~: $4$.

$Q_3$~: $\frac{3 \times 12}{4} = 9$, donc $Q_3$ est la $9$-ième
valeur~: $8$.

Diagramme en boîte~: boîte de $4$ à $8$ avec la barre de \omterm{def:g10:stats:median}{médiane} en $6$,
moustaches de $1$ à $11$.
\end{solution}

\begin{solution}{exo:g10:stats:4}
Les $24$ notes totalisent $24 \times 11.5 = 276$. Avec la nouvelle note
le total est $276 + 19 = 295$ pour $25$ élèves~: nouvelle \omterm{def:g10:stats:mean}{moyenne}
$\frac{295}{25} = 11.8$.
\end{solution}

\begin{solution}{exo:g10:stats:5}
Total actuel~: $5 \times 12 = 60$. Une \omterm{def:g10:stats:mean}{moyenne} de $13$ sur $6$ tests
exige un total de $78$, donc la $6$-ième note doit être $78 - 60 = 18$.
Une \omterm{def:g10:stats:mean}{moyenne} de $14$ exigerait $84 - 60 = 24 > 20$~: non atteignable.
\end{solution}

\begin{solution}{exo:g10:stats:6}
\omterm{def:g10:stats:mean}{Moyenne} plus grande que $20$~: par exemple
$1,\ 2,\ 3,\ 10,\ 40,\ 50,\ 60$ --- \omterm{def:g10:stats:median}{médiane} $10$ (la $4$-ième valeur
triée), \omterm{def:g10:stats:mean}{moyenne} $\frac{166}{7} \approx 23.7$.

\omterm{def:g10:stats:mean}{Moyenne} plus petite que $6$~: par exemple
$0,\ 0,\ 0,\ 10,\ 10,\ 10,\ 10$ --- \omterm{def:g10:stats:median}{médiane} $10$, \omterm{def:g10:stats:mean}{moyenne}
$\frac{40}{7} \approx 5.7$.

La \omterm{def:g10:stats:mean}{moyenne} ne peut pas descendre sous $\frac{40}{7}$~: pour que la
\omterm{def:g10:stats:median}{médiane} (la $4$-ième valeur triée) soit $10$, les valeurs en positions
$4$ à $7$ doivent toutes être $\geq 10$, donc le total est au moins
$40$ et la \omterm{def:g10:stats:mean}{moyenne} au moins $\frac{40}{7}$ --- le second exemple est
extrémal.

Ces exemples montrent que la \omterm{def:g10:stats:mean}{moyenne} et la \omterm{def:g10:stats:median}{médiane} sont largement
indépendantes~: connaître l'une dit peu sur l'autre, d'où l'intérêt d'un
bon résumé qui rapporte les deux.
\end{solution}

\begin{solution}{exo:g10:stats:7}
\emph{1.} \omterm{def:g10:stats:mean}{Moyenne}~:
$\frac{9 \times 1800 + 12000}{10} = \frac{16200 + 12000}{10} = 2820$.
\omterm{def:g10:stats:median}{Médiane}~: triés, les $10$ salaires sont neuf valeurs $1800$ puis
$12000$~; la \omterm{def:g10:stats:median}{médiane} est le \omterm{prop:g10:coordgeom:midpoint}{milieu} des $5$-ième et $6$-ième valeurs,
toutes deux $1800$~: \omterm{def:g10:stats:median}{médiane} $1800$.

\emph{2.} La \omterm{def:g10:stats:median}{médiane} ($1800$) décrit le salaire typique~: c'est ce que
gagnent réellement $9$ employés sur $10$. La \omterm{def:g10:stats:mean}{moyenne} ($2820$) est tirée
vers le haut par l'unique haut salaire et ne correspond à la fiche de
paie de personne.
\end{solution}

\begin{solution}{exo:g10:stats:8}
Taille totale~: $12 \times 168 + 18 \times 163 = 2016 + 2934 = 4950$~cm
pour $30$ élèves, donc la \omterm{def:g10:stats:mean}{moyenne} est $\frac{4950}{30} = 165$~cm. Elle
est plus proche de la \omterm{def:g10:stats:mean}{moyenne} des filles car elles sont plus nombreuses
(\omterm{def:g10:stats:mean}{moyenne} pondérée, poids $12$ et $18$).
\end{solution}

\begin{solution}{exo:g10:stats:9}
\emph{1.} La nouvelle \omterm{def:g10:stats:mean}{moyenne} est
\[
\bar y = \frac{y_1 + \dots + y_N}{N}
= \frac{(a x_1 + b) + \dots + (a x_N + b)}{N}
= \frac{a(x_1 + \dots + x_N) + Nb}{N}
= a \bar x + b .
\]

\emph{2.} Avec $a = 1.8$ et $b = 32$~: \omterm{def:g10:stats:mean}{moyenne}
$1.8 \times 20 + 32 = 68$ degrés Fahrenheit.
\end{solution}

\begin{solution}{pb:g10:stats:1}
\textbf{1.} \omterm{def:g10:stats:mean}{Moyenne}~: $\frac{2\,550}{7} \approx 364$ milliers
d'euros~; \omterm{def:g10:stats:median}{médiane}~: la quatrième des sept valeurs ordonnées, soit
$230$. C'est la \omterm{def:g10:stats:median}{médiane} qui décrit la rue --- six des sept maisons
sont très loin de $364$.

\textbf{2.} Sans le manoir~: \omterm{def:g10:stats:mean}{moyenne} $\frac{1\,350}{6} = 225$,
\omterm{def:g10:stats:median}{médiane} $\frac{220 + 230}{2} = 225$. La \omterm{def:g10:stats:mean}{moyenne} a chuté d'environ
$139$~; la \omterm{def:g10:stats:median}{médiane} de $5$. Une seule valeur extrême peut emmener la
\omterm{def:g10:stats:mean}{moyenne} n'importe où~; la \omterm{def:g10:stats:median}{médiane} s'en aperçoit à peine --- elle est
\emph{robuste}.

\textbf{3.} \omterm{def:g10:stats:mean}{Moyenne} $= \frac{0 \times 20 + 1 \times 25 + 2 \times
30 + 3 \times 15 + 4 \times 10}{100} = \frac{170}{100} = 1.7$.
Aucune famille n'a $1.7$ enfant, mais $1.7 \times 100 = 170$
\emph{est} le nombre total exact d'enfants~: la \omterm{def:g10:stats:mean}{moyenne} code le
total, redistribué également --- c'est là son vrai talent.

\textbf{4.} $N = 15$~: \omterm{def:g10:stats:median}{médiane} $=$ 8\textsuperscript{e} valeur
$= 11$~; $Q_1 =$ 4\textsuperscript{e} valeur $= 8$~;
$Q_3 =$ 12\textsuperscript{e} valeur $= 14$.

\textbf{5.} \omterm{def:g10:stats:quartiles}{Étendue} $19 - 4 = 15$~: la dispersion totale, extrêmes
compris. \omterm{def:g10:stats:quartiles}{Écart interquartile} $14 - 8 = 6$~: la largeur de la moitié
centrale --- la dispersion débarrassée du bruit des extrêmes.

\textbf{6.} Une translation de $+b$ déplace \emph{toutes} les
valeurs~; les différences de valeurs --- \omterm{def:g10:stats:quartiles}{étendue}, \omterm{def:g10:stats:quartiles}{écart
interquartile} --- sont donc inchangées~:
$(x_i + b) - (x_j + b) = x_i - x_j$. Une multiplication par $a$
multiplie toutes les différences par $a$~: \omterm{def:g10:stats:quartiles}{étendue} et \omterm{def:g10:stats:quartiles}{écart
interquartile} sont multipliés par $\abs a$.

\textbf{7.} Barème additif~: $6 \to 9$ et $18 \to 21$ --- l'élève
faible gagne $50\,\%$, le fort $17\,\%$, et $21$ ne crève que très
légèrement l'échelle sur $20$. Barème multiplicatif~: $6 \to 8$ et
$18 \to 24$ --- les rapports sont préservés, mais $24$ est
impossible sur une échelle de $20$~: c'est le problème du plafond.
Les barèmes additifs flattent le bas~; les multiplicatifs font
exploser le haut.

\textbf{8.} Total des notes~: $20 \times 12 + 30 \times 10 = 540$,
donc la \omterm{def:g10:stats:mean}{moyenne} globale vaut $\frac{540}{50} = 10.8$~: les trente
élèves de la classe B pèsent plus que les vingt de la classe A ---
les \omterm{def:g10:stats:mean}{moyennes} se regroupent par \omterm{def:g10:stats:mean}{moyenne} \emph{pondérée}, jamais par
\omterm{def:g10:stats:mean}{moyenne} simple.

\textbf{9.} \omterm{def:g10:stats:median}{Médiane} de $A$~: $2$~; de $B$~: $5$. Liste réunie
$\{0, 1, 2, 3, 10\}$~: \omterm{def:g10:stats:median}{médiane} $2$. La \omterm{def:g10:stats:mean}{moyenne} des \omterm{def:g10:stats:median}{médianes} pondérée
par les effectifs vaudrait
$\frac{3 \times 2 + 2 \times 5}{5} = 3.2 \neq 2$~: les \omterm{def:g10:stats:median}{médianes} ne
portent aucun total, aucune formule de regroupement n'existe donc ---
il faut retrier toute la liste.

\textbf{10.} \omterm{def:g10:stats:mean}{Moyenne}~: $\frac{31.3}{5} = 6.26$. \omterm{def:g10:stats:median}{Médiane}~: $5.4$.
\omterm{def:g10:stats:mean}{Moyenne} élaguée~: $\frac{5.3 + 5.4 + 5.5}{3} = 5.4$. Un juge
enthousiaste (ou corrompu) a déplacé la \omterm{def:g10:stats:mean}{moyenne} de près d'un point~;
la \omterm{def:g10:stats:median}{médiane} et la \omterm{def:g10:stats:mean}{moyenne} élaguée n'ont pas bronché --- d'où les
règles d'élagage des sports jugés.

\textbf{11.} Usine A~: moitié centrale à $0.2$ mm près~; usine B~: à
$4$ mm près. Même \omterm{def:g10:stats:mean}{moyenne}, fiabilité opposée~: acheter chez A. C'est
l'\omterm{def:g10:stats:quartiles}{écart interquartile} qui a tranché --- les \omterm{def:g10:stats:mean}{moyennes} seules comparent
des centres, jamais des régularités.

\textbf{12.} Une \omterm{def:g10:stats:median}{médiane} très inférieure à la \omterm{def:g10:stats:mean}{moyenne} trahit une
\emph{longue traîne à droite}~: la plupart des attentes sont courtes,
quelques-unes catastrophiques. Exemple~: $\{5, 5, 10, 10, 95\}$ ---
\omterm{def:g10:stats:median}{médiane} $10$, \omterm{def:g10:stats:mean}{moyenne} $\frac{125}{5} = 25$.

\textbf{13.} Cela signifie que $75\,\%$ des bébés du même âge sont
plus petits (et $25\,\%$ plus grands)~: c'est un rang, pas une
taille. «~Tous les enfants au-dessus de la \omterm{def:g10:stats:mean}{moyenne}~» est
impossible~: si toutes les valeurs dépassaient la \omterm{def:g10:stats:mean}{moyenne}, leur
\omterm{def:g10:stats:mean}{moyenne} dépasserait la \omterm{def:g10:stats:mean}{moyenne}, c'est-à-dire elle-même~:
contradiction. Mais \emph{presque} tous le peuvent~: un seul enfant
très petit suffit à maintenir la \omterm{def:g10:stats:mean}{moyenne} au-dessous de tous les
autres. Au-dessus de la \emph{\omterm{def:g10:stats:median}{médiane}}~: jamais plus de la moitié,
par définition. On ne flatte pas un rang~; on flatte une \omterm{def:g10:stats:mean}{moyenne}.

\textbf{14.} Note typique~: la même, \omterm{def:g10:stats:median}{médiane} $11$ pour les deux.
Homogénéité du centre~: B est plus resserrée (\omterm{def:g10:stats:quartiles}{écart interquartile} $3$
contre $5$). Extrêmes~: B est plus turbulente (\omterm{def:g10:stats:quartiles}{étendue} $17$ contre
$13$). B est une classe d'élèves semblables plus quelques valeurs
aberrantes~; A est régulièrement étalée.

\textbf{15.} (1) $5\,000$ est-il la \omterm{def:g10:stats:mean}{moyenne} ou la \omterm{def:g10:stats:median}{médiane} --- et
puis-je avoir la \omterm{def:g10:stats:median}{médiane}~? (2) Quelle est la dispersion (les
\omterm{def:g10:stats:quartiles}{quartiles}), pour savoir à quel point «~typique~» est typique~?
(3) Qui a été compté --- taille et sélection de l'échantillon (la
leçon sur les sondages du devoir du week-end sur les mensonges des
données, dans le volume précédent)~?

\textbf{16.} Par exemple $100, 150, 200, 250, 800$~: \omterm{def:g10:stats:median}{médiane} $200$,
\omterm{def:g10:stats:mean}{moyenne} $\frac{1\,500}{5} = 300$.

\textbf{17.} $\sum (x_i - \bar x) = \sum x_i - N \bar x =
N\bar x - N\bar x = 0$. Vérification~:
$-200 - 150 - 100 - 50 + 500 = 0$. La \omterm{def:g10:stats:mean}{moyenne} est le point où les
écarts s'équilibrent --- le centre de gravité des données (le $G$ du
devoir du week-end correspondant du volume précédent, à une
dimension).

\textbf{18.} À la \omterm{def:g10:stats:median}{médiane} $4$~: $3 + 1 + 0 + 5 + 9 = 18$ km. À la
\omterm{def:g10:stats:mean}{moyenne} $6$~: $5 + 3 + 2 + 3 + 7 = 20$ km. En $5$~:
$4 + 2 + 1 + 4 + 8 = 19$ km. La \omterm{def:g10:stats:median}{médiane} gagne --- s'en éloigner
revient toujours à s'écarter de plus de boutiques qu'on ne s'en
rapproche, si bien que le total ne peut que croître.

\textbf{19.} $0.3^{10} \approx 6 \times 10^{-6}$~: environ six
pilotes \emph{par million}. Sur $4\,000$ pilotes, l'effectif attendu
est $0.02$ --- c'est-à-dire zéro, comme l'a constaté le lieutenant
Daniels. L'armée de l'air a cessé de concevoir pour l'homme moyen et
a tout rendu réglable~: sièges, pédales, sangles --- concevoir pour
la \emph{dispersion}, non pour le centre.

\textbf{20.} Totaux~: la \emph{\omterm{def:g10:stats:mean}{moyenne}}. Individu typique~: la
\emph{\omterm{def:g10:stats:median}{médiane}}. Équité de la dispersion~: les \emph{\omterm{def:g10:stats:quartiles}{quartiles}} (et
l'\omterm{def:g10:stats:quartiles}{écart interquartile}). Contamination~: la \emph{\omterm{def:g10:stats:mean}{moyenne} élaguée}.
Et le mot de la fin~: l'individu moyen n'existe pas --- mais
l'individu médian, presque.
\end{solution}
